% English abstract before the 2026-11-08 TCS-style rewrite (README (215))
\begin{abstract}
We study how much positive data is needed to learn multiple context-free
languages when constituents are compared only if they agree under a typing
fixed in advance: a morphism onto a finite semigroup.  Identification in the
limit follows from Yoshinaka's distributional framework; we measure the total
length of characteristic samples against the size and rule thickness of a
grammar for the target, and ask for which typings this length is polynomial.
If the typing is onto a local group, a finite semigroup $\mathsf S$ whose local
monoids $e\mathsf Se$, for idempotents $e$, are groups, then for every fixed
fan-out and rank bound, samples of length quadratic in grammar size and linear
in rule thickness suffice; hence these classes are identifiable in polynomial
time and thick data.  The key step is a simultaneous shortening of interleaved
tuple components that preserves ordered products in noncommutative groups.
The broader typing regimes are witnessed by non-context-free targets of fan-out
two: one lies in the class of a local-group typing but of no window--group
typing, another in the class of a window--group typing over a noncommutative
group but of no window--abelian one.  A
factor linear in both grammar size and rule thickness is unavoidable under
every typing over large enough alphabets.  At fan-out one and rank at least
two, local groups are the exact boundary: a pseudovariety of finite semigroups
yields polynomial data for all its typings if and only if it consists of local
groups; otherwise some targets need samples exponentially longer than their
grammars, for every set-driven reconstruction operator with characteristic
samples throughout the class.  Whether this boundary persists at fan-out two
is open; the fan-out-one lower-bound family provably does not lift.
\end{abstract}

% Japanese abstract before the same rewrite
\begin{abstract}
本稿は，前もって固定した型付け（有限半群の上への準同型）の値が一致するときにのみ構成素を比較するという条件の下で，
多重文脈自由言語の学習にどれだけの正データが要るかを調べる．
極限同定は Yoshinaka の分布的枠組みから従う．
本稿は特徴標本の総長を標的の文法の大きさと規則の厚さに対して測り，どの型付けでこの長さが多項式になるかを問う．
型付けが局所群，すなわち有限半群 $\mathsf S$ であってその局所モノイド $e\mathsf Se$（$e$ は冪等元）がすべて群であるもの，の上への写像ならば，
固定したすべてのファンアウトとランクの上限について，文法の大きさの2乗と規則の厚さの1乗に比例する長さの標本で足り，
したがってこれらのクラスは多項式時間・厚データ同定可能である．
鍵となるのは，交互に並ぶタプル成分を，非可換群における順序付き積を保ったまま同時に短縮する操作である．
広い型付けの領域は，ファンアウト2の非文脈自由の標的で確かめられる：一つは局所群型付けのクラスには入るがどの窓--群型付けのクラスにも入らず，もう一つは非可換群上の窓--群型付けのクラスには入るがどの窓--可換群型付けのクラスにも入らない．
文法の大きさと規則の厚さの両方について線形な因子は，十分大きなアルファベット上ではどの型付けの下でも避けられない．
ファンアウト1かつランク2以上では，局所群がちょうど境界になる：
有限半群の擬多様体がそのすべての型付けで多項式のデータを与えるのは，それが局所群からなるとき，かつそのときに限る；
そうでなければ，クラス全体で特徴標本をもつすべての集合駆動型再構成作用素について，文法より指数的に長い標本を要する標的が存在する．
この境界がファンアウト2でも保たれるかは未解決であり，ファンアウト1の下界の族はそこへ持ち上がらないことが証明される．
\end{abstract}
